\chapter{Hito 6: Construcción del Primer Proceso ETL}

\section{Objetivo}

El objetivo de este hito es construir el primer flujo completo de carga de datos.

Por primera vez conectaremos:

\begin{verbatim}
CSV
|
V
Python
|
V
PostgreSQL
\end{verbatim}

Al finalizar este hito el estudiante será capaz de:

\begin{itemize}
\item Leer archivos CSV mediante Python.
\item Conectarse a PostgreSQL utilizando SQLAlchemy.
\item Cargar información en la capa staging.
\item Automatizar el proceso de carga.
\item Validar que los datos fueron cargados correctamente.
\end{itemize}

\section{Resultado esperado}

Al finalizar este hito los datos contenidos en:

\begin{verbatim}
data/raw/tipo_cambio.csv

data/raw/balances.csv
\end{verbatim}

deberán encontrarse cargados en:

\begin{verbatim}
staging.tipo_cambio_raw

staging.balance_bancario_raw
\end{verbatim}

\section{Arquitectura del proceso}

El flujo que construiremos será:

\begin{verbatim}
CSV
|
V
Pandas
|
V
SQLAlchemy
|
V
PostgreSQL
\end{verbatim}

\section{Paso 1: Verificar archivos fuente}

Confirmar la existencia de:

\begin{verbatim}
data/raw/tipo_cambio.csv

data/raw/balances.csv
\end{verbatim}

Verificar que pueden abrirse correctamente.

\section{Paso 2: Crear módulo de conexión}

Verificar el archivo:

\begin{verbatim}
etl/db.py
\end{verbatim}

Contenido:

\begin{lstlisting}[language=Python]
from sqlalchemy import create_engine
from dotenv import load_dotenv

import os

load_dotenv()

HOST = os.getenv("DB_HOST")
PORT = os.getenv("DB_PORT")
DB = os.getenv("DB_NAME")
USER = os.getenv("DB_USER")
PASSWORD = os.getenv("DB_PASSWORD")

connection_string = (
f"postgresql+psycopg2://"
f"{USER}:{PASSWORD}@{HOST}:{PORT}/{DB}"
)

engine = create_engine(
connection_string
)
\end{lstlisting}

\section{Paso 3: Crear script de carga}

Crear:

\begin{verbatim}
etl/load_staging.py
\end{verbatim}

\section{Paso 4: Importar librerías}

Agregar:

\begin{lstlisting}[language=Python]
import pandas as pd

from db import engine
\end{lstlisting}

\section{Paso 5: Leer tipo de cambio}

Agregar:

\begin{lstlisting}[language=Python]
fx = pd.read_csv(
"data/raw/tipo_cambio.csv"
)

print(
f"Filas tipo cambio: {len(fx)}"
)
\end{lstlisting}

\section{Paso 6: Leer balances}

Agregar:

\begin{lstlisting}[language=Python]
balances = pd.read_csv(
"data/raw/balances.csv"
)

print(
f"Filas balances: {len(balances)}"
)
\end{lstlisting}

\section{Paso 7: Agregar metadatos}

Agregar:

\begin{lstlisting}[language=Python]
fx["archivo_origen"] = (
"tipo_cambio.csv"
)

balances["archivo_origen"] = (
"balances.csv"
)
\end{lstlisting}

Estos metadatos permiten rastrear el origen de cada carga.

\section{Paso 8: Limpiar staging antes de cargar}

Crear:

\begin{verbatim}
sql/truncate_staging.sql
\end{verbatim}

Contenido:

\begin{lstlisting}[language=SQL]
TRUNCATE TABLE
staging.tipo_cambio_raw;

TRUNCATE TABLE
staging.balance_bancario_raw;
\end{lstlisting}

Esta estrategia garantiza idempotencia.

\section{Paso 9: Ejecutar truncado desde Python}

Agregar:

\begin{lstlisting}[language=Python]
from sqlalchemy import text

with engine.begin() as conn:

```
conn.execute(
    text(
        """
        TRUNCATE TABLE
        staging.tipo_cambio_raw;

        TRUNCATE TABLE
        staging.balance_bancario_raw;
        """
    )
)
```

\end{lstlisting}

\section{Paso 10: Cargar tipo de cambio}

Agregar:

\begin{lstlisting}[language=Python]
fx.to_sql(
"tipo_cambio_raw",
con=engine,
schema="staging",
if_exists="append",
index=False
)
\end{lstlisting}

\section{Paso 11: Cargar balances}

Agregar:

\begin{lstlisting}[language=Python]
balances.to_sql(
"balance_bancario_raw",
con=engine,
schema="staging",
if_exists="append",
index=False
)
\end{lstlisting}

\section{Paso 12: Confirmar carga}

Agregar:

\begin{lstlisting}[language=Python]
print(
"Carga completada."
)
\end{lstlisting}

\section{Versión completa del script}

El archivo completo debe verse así:

\begin{lstlisting}[language=Python]
import pandas as pd

from sqlalchemy import text

from db import engine

fx = pd.read_csv(
"data/raw/tipo_cambio.csv"
)

balances = pd.read_csv(
"data/raw/balances.csv"
)

fx["archivo_origen"] = (
"tipo_cambio.csv"
)

balances["archivo_origen"] = (
"balances.csv"
)

with engine.begin() as conn:

```
conn.execute(
    text(
        """
        TRUNCATE TABLE
        staging.tipo_cambio_raw;

        TRUNCATE TABLE
        staging.balance_bancario_raw;
        """
    )
)
```

fx.to_sql(
"tipo_cambio_raw",
con=engine,
schema="staging",
if_exists="append",
index=False
)

balances.to_sql(
"balance_bancario_raw",
con=engine,
schema="staging",
if_exists="append",
index=False
)

print("Carga completada.")
\end{lstlisting}

\section{Paso 13: Ejecutar ETL}

Desde terminal:

\begin{lstlisting}[language=bash]
python etl/load_staging.py
\end{lstlisting}

Salida esperada:

\begin{verbatim}
Carga completada.
\end{verbatim}

\section{Paso 14: Verificar en PostgreSQL}

Consultar:

\begin{lstlisting}[language=SQL]
SELECT COUNT(*)
FROM staging.tipo_cambio_raw;
\end{lstlisting}

Consultar:

\begin{lstlisting}[language=SQL]
SELECT COUNT(*)
FROM staging.balance_bancario_raw;
\end{lstlisting}

Los conteos deben coincidir con los CSV.

\section{Paso 15: Inspeccionar datos}

Consultar:

\begin{lstlisting}[language=SQL]
SELECT *
FROM staging.tipo_cambio_raw
LIMIT 10;
\end{lstlisting}

Consultar:

\begin{lstlisting}[language=SQL]
SELECT *
FROM staging.balance_bancario_raw
LIMIT 10;
\end{lstlisting}

\section{Paso 16: Crear script de validación}

Crear:

\begin{verbatim}
sql/check_load.sql
\end{verbatim}

Contenido:

\begin{lstlisting}[language=SQL]
SELECT
COUNT(*) AS total_fx
FROM staging.tipo_cambio_raw;

SELECT
COUNT(*) AS total_balances
FROM staging.balance_bancario_raw;
\end{lstlisting}

\section{Paso 17: Crear orquestador}

Crear:

\begin{verbatim}
etl/run_pipeline.py
\end{verbatim}

Contenido inicial:

\begin{lstlisting}[language=Python]
import subprocess

subprocess.run(
[
"python",
"etl/load_staging.py"
]
)
\end{lstlisting}

Este archivo será el punto de entrada de todo el pipeline.

\section{Buenas prácticas}

Durante este proyecto seguiremos las siguientes reglas:

\begin{enumerate}
\item El ETL debe ser reproducible.
\item La carga debe ser idempotente.
\item No modificar datos manualmente.
\item Validar cada carga.
\item Mantener trazabilidad mediante metadatos.
\end{enumerate}

\section{Checklist de validación}

Antes de continuar verificar:

\begin{itemize}
\item[$\Box$] Python puede leer ambos CSV.
\item[$\Box$] Python puede conectarse a PostgreSQL.
\item[$\Box$] Staging se limpia correctamente.
\item[$\Box$] Los datos son cargados.
\item[$\Box$] Los conteos coinciden.
\item[$\Box$] El script puede ejecutarse múltiples veces.
\end{itemize}

\section{Entregable}

Al finalizar este hito el estudiante deberá disponer de:

\begin{itemize}
\item Script de carga funcional.
\item Datos cargados en staging.
\item Validaciones básicas.
\item Pipeline reproducible.
\item Primer flujo ETL completo.
\end{itemize}

\section{Próximo hito}

En el Hito 7 construiremos la capa Core.

Crearemos:

\begin{itemize}
\item Dimensión Tiempo.
\item Dimensión Moneda.
\item Dimensión Entidad Financiera.
\item Claves subrogadas.
\item Primera transformación desde staging hacia el modelo dimensional.
\end{itemize}
